% !TEX root = ../ApproxDA-TransUNet.tex
\section{Conclusion}
\label{sec:conclusion}
This paper presented ApproxDA-TransUNet, an efficient dual-attention approximation framework for medical image segmentation. Across three representative benchmarks, ApproxDA-TransUNet consistently outperformed the original DA-TransUNet while retaining efficient attention computation, demonstrating that appropriately designed attention approximation can improve segmentation accuracy rather than merely reducing computational cost.
Beyond the proposed architecture, our experiments provide new insights into attention approximation. We show that approximation effectiveness depends on the contextual requirements of the target task, which we describe as Global Context Sensitivity (GCS). High-GCS tasks require careful approximation design, whereas low-GCS tasks remain robust across a broad range of approximation scales. We further demonstrate that conventional scalar gate learning collapses toward nearly uniform routing because of gradient symmetry, suggesting that effective adaptive fusion requires mechanisms that explicitly promote branch specialization.
Beyond the core results, this journal extension presents two additional analyses. A training dynamics study (Section~\ref{sec:generalization}) shows that windowed attention acts as a task-dependent regularizer: on low-GCS tasks it reduces val DSC volatility by 9--18$\times$ relative to global attention, while slightly increasing volatility on high-GCS Synapse. This CS-dependent stability crossover is consistent with the inductive bias interpretation and holds across all gate=pam window configurations. A causal elimination study (Section~\ref{sec:gcs_mechanism}) systematically rules out image texture, object size, spatial dispersion, and component-level window crossing as GCS drivers, converging on Semantic Spatial Dependency (SSD)---the requirement to locate each semantic structure relative to others---as the unifying mechanism. A discriminating experiment on ACDC cardiac MRI (4 classes, spatially co-located anatomy) yields $\Delta\text{DSC}{=}0.73$\,pp, clustered with binary tasks ($0.50$--$0.64$\,pp) and far below Synapse ($2.30$\,pp), directly confirming that spatial spread of classes---not class count---governs GCS.

Our results suggest that approximation can actively improve representation learning by introducing a task-aligned locality prior---shifting its role from an engineering optimization to a model-design principle. The primary open direction is extending the GCS spectrum to tasks with intermediate spatial dependency structures (e.g., multi-class datasets with partially distributed anatomy), and designing attention routing mechanisms that explicitly promote branch diversity to avoid gradient symmetry collapse. Future research should first ask \emph{how much context} a task actually requires---a question GCS and SSD can answer from annotation structure alone---before deciding how to approximate.

\textbf{Limitations.}
Three limitations of the current work warrant attention.
First, the locality prior introduced by windowed attention benefits locally-resolvable structures but can degrade those requiring global spatial extent estimation. On ACDC cardiac MRI (GCS~$=$~0.73\,pp), ApproxDA-TransUNet achieves the best per-class scores on RV (\textbf{89.61\%}) and Myo (\textbf{85.91\%}), whose boundaries are resolved by local intensity gradients, while the LV cavity score (91.40\%) falls 3--4\,pp below competitive baselines---likely because LV delineation requires estimating global cavity extent across multiple attention windows. This intra-dataset heterogeneity shows that GCS is not monolithic: individual structures within a low-GCS dataset may differ in sub-structure context sensitivity. Future work should measure GCS at the per-class level to distinguish locally-bounded from globally-defined structures within the same scan.
Second, the low-rank projection ($r{=}32$) may degrade attention fidelity for thin or fine-grained structures such as the myocardium, where subpixel boundary precision matters more than context scope. Adaptive rank selection---higher $r$ for precision-critical tasks---is a natural extension.
Third, measuring GCS empirically still requires running a window-size ablation (four training runs per dataset), which is not zero-cost. Predicting GCS directly from mask statistics---using SSD proxies such as SC5---is an open problem that would make the framework fully annotation-driven without additional training.